\documentclass[10pt,a4paper]{amsart}
\usepackage[margin=19mm]{geometry}
\usepackage{amsmath,amssymb,mathtools}
\usepackage[scr=rsfs,cal=euler]{mathalfa}
\usepackage{xfrac}
\usepackage[unicode,hidelinks,bookmarks=false]{hyperref}
\newcommand{\ie}{i.e.\ }
\newcommand{\cf}{cf.\ }
\newcommand{\resp}{respectively\ }
\newcommand{\Subsection}{\subsection}
\providecommand{\nobreakdash}{\nobreak\hbox{-}}
\setlength{\emergencystretch}{2em}
\multlinegap0pt
\usepackage{bm}
\usepackage{bbm}
\usepackage{dsfont}
\usepackage{xspace}
\usepackage{cancel}
\usepackage{mathtools}
\usepackage{amsthm}
\makeatletter
\@ifundefined{theorem}{\newtheorem{theorem}{Theorem}}{}
\@ifundefined{lemma}{\newtheorem{lemma}[theorem]{Lemma}}{}
\@ifundefined{proposition}{\newtheorem{proposition}[theorem]{Proposition}}{}
\makeatother
\newcommand{\R}{\mathbb R}
\title[Geometric properties of solutions to elliptic PDE's in Gauss]{Geometric properties of solutions to elliptic PDE's in Gauss space and related Brunn-Minkowski type inequalities}
\author{Andrea Colesanti, Lei Qin, Paolo Salani}
\date{Source: arXiv:2502.00184v2 (CC BY 4.0)}
\begin{document}
\maketitle

\noindent\emph{Source and licence.} Geometric properties of solutions to elliptic PDE's in Gauss space and related Brunn-Minkowski type inequalities. Version arXiv:2502.00184v2 (CC BY 4.0), distributed under the Creative Commons Attribution 4.0 International licence, as recorded in the arXiv licence metadata for this exact version and shown on the abstract page https://arxiv.org/abs/2502.00184v2. Full rights notice: licenses/arxiv-2502.00184-CC-BY-4.0.txt.\par\smallskip

\noindent\textit{Editorial scope.} Counted unit: the Lemma whose source label is \texttt{Caccioppoli} in the article (source0.tex lines 542--552, source label \texttt{Caccioppoli}), delivered together with the article's own complete original proof of it (source0.tex lines 554--590, one \texttt{proof} environment of 37 lines). No mathematics is written, repaired or re-derived: statement and proof are the article's own text line for line. Required context retained uncounted: PDE1 (lines 130--137); inf-convo-prop (lines 387--424); Viscosity1 (lines 430--439). Every definition, display and label of those lines is retained unchanged. Omitted from this package and not redistributed: 1--129, 138--386, 425--429, 440--541, 553--553, 591--1146. Nothing inside a retained excerpt is omitted or altered: every retained body is delivered exactly as the article's lines are written, so no omission receipt of any kind accompanies this package. Citations: the retained excerpts carry no citation command at all, so no bibliography entry of the article is transcribed and no reference is redistributed. Reference adaptation: none was needed and none was made; every reference inside the delivered unit resolves inside it, so no unresolved reference or citation is produced. Printed slips kept literal: no printed word, symbol, hypothesis or number of the counted statement or of its proof was altered. Layout and class: the article's own class is \texttt{amsart} with packages including geometry, amsmath, amssymb, amscd, amsthm, latexsym, graphicx, babel, inputenc, textcomp, times, colortbl; the wrapper is the lane's pinned \texttt{amsart} editorial wrapper at 10pt on A4, which restates the theorem-like environments the retained text needs and adds the article's own macro definitions that the retained text needs (\texttt{\textbackslash R}), with extra wrapper packages (bm, bbm, dsfont, xspace, cancel, mathtools). The delivered numbering is the unit's own. Assets: no retained span contains a figure environment, a graphics-inclusion command, a PDF-page inclusion, a table or a TikZ picture; no image file is copied into this package, the package declares no asset dependency and no image file is redistributed with it. Licence: version arXiv:2502.00184v2 of this article is distributed under the Creative Commons Attribution 4.0 International licence, as recorded in the arXiv licence metadata for this exact version and shown on the abstract page https://arxiv.org/abs/2502.00184v2. A full rights notice is delivered at licenses/arxiv-2502.00184-CC-BY-4.0.txt. Characters: every retained excerpt is pure ASCII, so the delivered engine, XeLaTeX with its default Latin Modern text font, reports no missing character.
\begin{equation}\label{PDE1}
\left\{
\begin{aligned}
& -\Delta_{p,\gamma} u=\lambda_{p,\gamma} |u|^{p-2}u \ \  &\mathrm {in}\ \ \Omega, \\
&\ u=0 \ \ \ \ &\mathrm {on} \ \ \partial \Omega\,,
\end{aligned}
\right.
\end{equation}

\begin{proposition}\label{inf-convo-prop}
Suppose that $u:\Omega \rightarrow \R$ is a bounded and upper semi-continuous function. Then the sup-convolution $u_\varepsilon$ satisfies the following properties:
\begin{itemize}
  \item The family $\{ u_\varepsilon\}$ is decreasing with respect to $\varepsilon$, $u_\varepsilon \geq u$ in $\Omega$ and $u_\varepsilon \rightarrow u$ locally uniformly as $\varepsilon \rightarrow  0$.

   \item $u_\varepsilon$ is locally Lipschitz and a.e. twice differentiable in the following sense: for almost every $x,y\in\Omega$,
      \[
      u_\varepsilon(y)=u_\varepsilon(x)+\nabla u_\varepsilon(x)\cdot (x-y)+\frac{1}{2}D^2 u_\varepsilon (x)(x-y)^2+o(|x-y|^2).
      \]

  \item There exists $r(\varepsilon )>0$ such that
  \[
  u_\varepsilon(x)=\mathop{\sup}_{y\in B_{r(\varepsilon)}(x)\cap \Omega} \left( u(y)-\frac{|x-y|^q}{q\varepsilon^{q-1}}\right), \]
  with $r(\varepsilon)\rightarrow 0 $ as $\varepsilon \rightarrow 0$.

  \item If $x\in \Omega_{r(\varepsilon)}:=\{ x\in \Omega: \mathrm{dist}(x,\partial \Omega)>r(\varepsilon)\}$, then there exists $x_\varepsilon\in B_{r(\varepsilon)}(x)$ such that
  \[
  u_{\varepsilon}(x)=u(x_\varepsilon)-\frac{1}{q\varepsilon^{q-1}}|x-x_\varepsilon|^q.
  \]
  Moreover, if the gradient $\nabla u_\varepsilon(x)$ exists, it holds
  \[
  \left( \frac{|x-x_\varepsilon|}{\varepsilon}\right)^{q-1} \leq |\nabla u_\varepsilon(x)|.
  \]
  In particular, if $\nabla u_\varepsilon (x)=0$, then $u_\varepsilon (x)=u(x)$.


  \item If $(\eta,X)\in J^{2,+} u_\varepsilon(x)$ with $x\in \Omega_{r(\varepsilon)}$, then
  \[
  \eta=\frac{|x-x_\varepsilon|^{q-2}(x-x_\varepsilon)}{\varepsilon^{q-1}},\ \mathop{and}\ X \geq - \frac{q-1}{\varepsilon}|\eta|^{\frac{q-2}{q-1}} I;
  \]
 here $I$ is the identity matrix of order $n$.
  \item $u_\varepsilon $ is semi-convex in $\Omega_{r(\varepsilon)}$, that is, there is a positive constant $c$ such that the function
  \[
  x\mapsto u_\varepsilon (x)+c|x|^2
  \]
  is convex, where the constant $c$ depends on $\varepsilon,\,q$ and the oscillation $\sup_{\Omega}-\inf_{\Omega} u$. Moreover, for a.e. $x\in \Omega_{r(\varepsilon)}$, $D^2 u_\varepsilon(x) \geq -cI$.
\end{itemize}
\end{proposition}

\begin{lemma}\label{Viscosity1}
Given $p>1$. Suppose that $u:\Omega \rightarrow \R$ is bounded, non-negative and upper semi-continuous function. Then, if $u$ is a viscosity sub-solution in $\Omega$ of the equation \eqref{PDE1}, then $u_\varepsilon$ is a viscosity sub-solution in $ \Omega_{r(\varepsilon)}$ of
\begin{equation}\label{viscosity-super-solution}
 -\Delta_p u +(x,\nabla u)|\nabla u|^{p-2}= f_\varepsilon(x,u,\nabla u),
\end{equation}
where
\begin{equation}\label{3-1}
f_\varepsilon(x,u_\varepsilon,\nabla u_\varepsilon):=\varepsilon |\nabla u_\varepsilon|^{p}+\lambda_{p,\gamma}\cdot{\sup}_{y\in B_{r(\varepsilon)}(x)} u_\varepsilon(y)^{p-1}.
\end{equation}
\end{lemma}

\begin{lemma}\label{Caccioppoli}
Under the same assumptions of Lemma \ref{Viscosity1}, {   If, for every nonnegative test function $\varphi\in C^{\infty}_0(\Omega)$, $u_\varepsilon$ satisfies the following inequality
\begin{equation}\label{newassumption}
\int_{\Omega}|\nabla u_\varepsilon|^{p-2}\nabla u_\varepsilon\cdot \nabla \varphi d\gamma \leq \int_\Omega  f_\varepsilon(x,u_\varepsilon,\nabla u_\varepsilon)\varphi d\gamma \,,
\end{equation}}
then, there exists a positive constant $c$, depending on $p,\,\lambda_{p,\,\gamma},\|u\|_{L^{\infty}(\Omega)},|\Omega|$,  such that for every test function $\psi\in C^{\infty}_0(\Omega_{r(\varepsilon)})$, $ 0\leq \psi \leq 1$, we have
\begin{equation}\label{Caccioppoli ineq}
\int_{\Omega} |\nabla u_\varepsilon|^p \psi^p d\gamma \leq c(1+\|\nabla \psi\|_{L^{\infty}(\Omega)}),
\end{equation}
for sufficiently small $\varepsilon>0$.
\end{lemma}

\begin{proof}[\bf Proof.]
Let $\psi \in C^{\infty}_0(\Omega_{r(\varepsilon)})$, $0\leq\psi\leq1$ and consider the test function
\[
\varphi(x):=(u_\varepsilon(x)-{\inf}_{K} u_\varepsilon  )\psi^p(x),\quad x\in \Omega.
\]
%{\color{blue} \{Delete it : Since $u_\varepsilon$ is twice differentiable a.e. in $\Omega$ and the pair $(Du_\varepsilon,D^2u_\varepsilon)$ belongs to the second "jets" $J^{2,+}u(x)$ and $J^{2,-}u(x)$ at the points of twice differentiability, after an integration by parts \} }
By assumption, we have
\[
\int_{\Omega}|\nabla u_\varepsilon|^{p-2}\nabla u_\varepsilon\cdot \nabla \varphi d\gamma \leq \int_\Omega  f_\varepsilon(x,u_\varepsilon,\nabla u_\varepsilon)\varphi d\gamma .
\]
By a direct calculation, the integral on the left hand side equals to
\begin{equation}\label{333}
\int_{\Omega} |\nabla u_\varepsilon|^{p-2}\nabla u_\varepsilon \cdot[ \psi^p \nabla u_\varepsilon+p\psi^{p-1}\nabla \psi ( u_\varepsilon(x)-{\inf}_{K} u_\varepsilon)]d\gamma.
\end{equation}
The Young's inequality
\[
ab \leq \delta a^q+\delta^{-1/(q-1)} b^{q'}
\]
where $\delta>0$, $q$ and $q'$ are conjugate exponents such that $1/q+1/q'=1$, implies that the absolute value of second term in \eqref{333} is bounded by
\[
\delta \int_\Omega |\nabla u_\varepsilon|^p \psi^p +\delta^{1-p} \int_\Omega p^p |\nabla \psi|^p (\mathrm{osc}_K u_\varepsilon)^pd\gamma.
\]
Therefore, we have
\begin{align}\label{3-6}
(1-\delta)\int_\Omega |\nabla u_\varepsilon|^p \psi^pd\gamma  \leq  \delta^{1-p}\int_\Omega p^p |\nabla \psi|^p (\mathrm{osc}_K u_\varepsilon)^pd\gamma+\int_\Omega  f_\varepsilon(x,u_\varepsilon,\nabla u_\varepsilon)\varphi d\gamma.
\end{align}
Observe that $\max\{|u_\varepsilon|, \mathrm{osc}_K u_\varepsilon\}\leq \|u\|_{L^{\infty}(\Omega)}$. By Proposition \ref{inf-convo-prop}, we have
\begin{align*}
\int_\Omega  f_\varepsilon(x,u_\varepsilon,\nabla u_\varepsilon)\varphi d\gamma&=\int_\Omega  (\varepsilon |\nabla u_\varepsilon|^p+\lambda_{p,\gamma}\cdot{\sup}_{y\in B_{r(\varepsilon)}(x)} u_\varepsilon(y)^{p-1} )\varphi\,d\gamma\\
& \leq \varepsilon\|u\|_{L^{\infty}(\Omega)}\int_\Omega |\nabla u_\varepsilon|^p  \psi^p  d\gamma+\lambda_{p,\gamma}\|u\|^p_{L^{\infty}(\Omega)}\cdot |\gamma(\Omega)|.
\end{align*}
Selecting $\delta <1/3$ and $\varepsilon\leq \frac{1}{3}\|u\|^{-1}_{L^{\infty}(\Omega)}$, by \eqref{3-6}, we have
\[
\int_\Omega |\nabla u_\varepsilon|^p \psi^pd\gamma  \leq c(1+\|\nabla \psi\|_{L^{\infty}(\Omega)}),
\]
where $c$ depends on $p,\lambda_{p,\gamma},\|u\|_{L^{\infty}(\Omega)},|\Omega|$, as desired.
\end{proof}

\end{document}
